\documentclass[10pt,a4paper]{amsart}
\usepackage[margin=19mm]{geometry}
\usepackage{amsmath,amssymb,mathtools}
\usepackage[scr=rsfs,cal=euler]{mathalfa}
\usepackage{xfrac}
\usepackage[unicode,hidelinks,bookmarks=false]{hyperref}
\newcommand{\ie}{i.e.\ }
\newcommand{\cf}{cf.\ }
\newcommand{\resp}{respectively\ }
\newcommand{\Subsection}{\subsection}
\providecommand{\nobreakdash}{\nobreak\hbox{-}}
\setlength{\emergencystretch}{2em}
\multlinegap0pt
\usepackage{amssymb}
\newtheorem{lem}{Lemma}
\newtheorem{theorem}{Theorem}
\newcommand\Aff{\mathfrak A}
\newcommand\BB{\textit{\textbf{B}}}
\newcommand\DD{\mathbb D}
\newcommand\Gb{\mathbf{G}}
\newcommand\Mbf{\mathbf M}
\newcommand\T{\mathcal T}
\newcommand\bal{\boldsymbol{\alpha}}
\newcommand\bep{\boldsymbol{\epsilon}}
\newcommand\cb{\mathbf{c}}
\newcommand\gb{\mathbf{g}}
\newcommand\lt{\mathfrak{l}}
\newcommand\nn{N}
\newcommand\rt{\mathfrak{r}}
\newcommand\ut{\mathfrak{u}}
\newcommand\vt{\mathfrak{v}}
\title{Notes on hyperbolic branching Brownian motion}
\author{Wolfgang Woess}
\date{Source: arXiv:2405.07301v5 (CC BY 4.0)}
\begin{document}
\maketitle

\noindent\emph{Source and licence.} Notes on hyperbolic branching Brownian motion. Version arXiv:2405.07301v5 (CC BY 4.0), distributed by arXiv under the Creative Commons Attribution 4.0 International licence, http://creativecommons.org/licenses/by/4.0/, as recorded in the arXiv licence metadata for this exact version and shown on the abstract page https://arxiv.org/abs/2405.07301v5. Full licence notice: \texttt{licenses/arxiv-2405.07301-CC-BY-4.0.txt}.\par\smallskip

\noindent\textit{Editorial scope.} Counted unit: the article's theorem thm:atoms (source0.tex lines 1759--1763). The article's complete original proof of it is delivered with it (source0.tex lines 1765--1823, one \texttt{proof} environment of 59 lines). No mathematics is written, repaired or re-derived: the statement and the proof are the article's own lines, in the article's own notation, and no retained line is substituted or re-flowed.  Retained uncounted as the source-local context the proof needs: lines 659--659; lines 741--744; lines 1713--1721; lines 1732--1737.  Nothing was omitted from any retained span, so the package carries no omission record.  No retained line contains a citation, so the wrapper carries no bibliography and no bibliography entry.  No retained line contains a figure environment, an image-inclusion command or drawing code, so no image is delivered and the package declares no asset dependency.  Editorial formatting only, in addition: an AMS \texttt{amsart} preamble that restates the article's own declarations and macros used by the retained lines, namely the environments \texttt{lem}, \texttt{theorem} on separate counters, together with the standard packages those lines need.  The retained environments print in this document as Lemma 1, Theorem 1 in order of appearance, whereas the article prints the same items with the numbering of its own full document.  The complete native source of the article as distributed in the arXiv e-print for this exact version is bundled unchanged as \texttt{sources/159033/source0.tex}.
\section{Hyperbolic branching Brownian motion}\label{sec:hypBBM}

\begin{equation}\label{eq:Gv}
%X_{\vt} 
\Gb_{\vt} =\gb_{\vt_1}\cdots\gb_{\vt_{k}}
\end{equation}

\begin{equation}\label{eq:muu}
%\begin{aligned}
\mu_{\ut, \infty}^{z} %&
= \lim_{t \to \infty} \mu_{\ut, t}^{z}\,, \quad \text{where }\; z=z_{\ut'} = \BB_{\ut'}\;\text{ and }\; %\\
\mu_{\ut, t}^{z} %&
= \frac{1}{\nn_{\ut}(t)} \sum_{\tau \in \T_{\ut}\,:\, |\tau| = t+ |u^-|} 
\delta_{\BB_{\tau}}
%\end{aligned}
\end{equation}

\begin{lem}\label{lem:indep}
Let $\ut, \vt \in  \{\lt,\rt\}^*$ be such that none of the two is an ancestor (predecessor) 
of the other.
Consider hyperbolic BBM indexed by the time trees $\T_{\ut}$ and $\T_{\vt}\,$, respectively, and starting at $o$. Then with probability $1$, the limit distributions 
$\mu_{\ut, \infty}^{o}$ and $\mu_{\vt, \infty}^{o}$ of \eqref{eq:muu} share no atoms.
\end{lem}

\begin{theorem}\label{thm:atoms}
With probability $1$, the support of the random limit distribution 
$\mu_{\infty}^{o} = \mu_{\bep, \infty}^{o}$
is infinite.
\end{theorem}

\begin{proof} For any measure $\mu$ on $\partial \DD$, we shall write 
$\mathsf{atoms}(\mu)$ for the set of atoms of  $\mu$, and $\mathsf{supp}(\mu)$
for its support.

Let $\ut \in  \{\lt,\rt\}^*$ and consider hyperbolic BBM according to our construction
of \S \ref{sec:hypBBM} with time tree $\T_{\ut}$ and starting point $\BB_{\ut'} = o$. 
The two children of $\ut$ are $\ut\lt$ and $\ut\rt$. For each of the 
two independent subtrees $\T_{\ut\lt}$ and $\T_{\ut\rt}\,$, the vertex
$\ut$ has the role which the ancestor $\bal$ has in $\T$. Write 
$z_1 = \BB_{\ut}\,$, so that $g_{z_1} = \gb_{\ut}\,$, the random element
of $\Aff$ as described in \S \ref{sec:hypBBM}.
Denote by $\Mbf_{\ut\lt,t}^{z_1}$ and $\Mbf_{\ut\rt,t}^{z_1}$ the occupation measures
at time $t$ (that is, distance $t$ from $u$ in $\T_{\lt}$ and $\T_{\rt}$),
respectively. 
Then
$$
\Mbf_{\ut,t}^o = \Mbf_{\ut\lt,t-\ell_{\ut}}^{z_1} + \Mbf_{\ut\rt,t-\ell_{\ut}}^{z_1}\,.
$$
Dividing by $\nn_{\bep}(t)$ and letting $t \to \infty\,$,
$$
\mu_{\ut,\infty}^o 
= \cb_{\ut\lt} \,\mu_{\ut\lt,\infty}^{z_1}
+ \cb_{\ut\rt}\,\mu_{\ut\rt,\infty}^{z_1},
\quad \text{where}\quad \cb_{\ut\lt} = \frac{W_{\ut\lt}}{W_{\ut}} \, e^{-\lambda\ell_{\ut}}
\;\text{ and }\; \cb_{\ut\rt} = \frac{W_{\ut\rt}}{W_{\ut}} \, e^{-\lambda\ell_{\ut}}\,,
$$
a non-trivial convex combination of two probability measures. 
We can also rewrite this as
$$
\gb_{\ut}^{-1}\mu_{\ut,\infty}^o 
= \cb_{\ut\lt} \,\mu_{\ut\lt,\infty}^{o}
+ \cb_{\ut\rt}\,\mu_{\ut\rt,\infty}^{o}\,.
$$
From this identity and Lemma \ref{lem:indep} we infer that
$$
\mathsf{atoms}(\mu_{\ut,\infty}^o) 
= \gb_{\ut}\mathsf{atoms}(\mu_{\ut\lt,\infty}^o)  \stackrel{{\,}_{_{_{_+}}}}{\cup} 
%\mathbin{\dot{\cup}}
\gb_{\ut}\mathsf{atoms}(\mu_{\ut\rt,\infty}^o)
$$
is almost surely a disjoint union. 

Recursively, we get the following convex combination for each $n$:
\begin{equation}\label{eq:meassum}
\mu_{\bep,\infty}^o = 
\sum_{\vt \in \{\lt,\rt\}^n} \cb^{(n)}_{\vt}\, \Gb_{\vt'}\, \mu_{\vt,\infty}^o
\end{equation}
with positive random constants $\cb^{(n)}_{\vt}$ and $\Gb_{v'}$ as defined by \eqref{eq:Gv}, and 
$$
\mathsf{atoms}(\mu_{\ut,\infty}^o) = %\mathbin{\dot{\bigcup}}
\ {\bigcup}_{\vt \in \{\lt,\rt\}^n}^{\!\!\!\!\!\!\!+}
%\stackrel{{\,\,}_{_{_{_{_{_+}}}}}}{\bigcup}_{\vt \in \{\lt,\rt\}^n}  
\Gb_{\vt'}\, \mathsf{atoms}(\mu_{\vt,\infty}^o) \quad \text{almost surely.}
$$
If $\mathsf{supp}(\mu_{\ut,\infty}^o)$ is finite then the support of the measure coincides with
the set of its atoms, and there must be $\vt$ such that $\mu_{\vt,\infty}^o$ has no atoms.
But along with $\mu_{\bep,\infty}^o\,$, by \eqref{eq:meassum} all $\mu_{\vt,\infty}^o$
have finite support consisting only of atoms, a contradiction.
\end{proof}

\end{document}
